%! AP Statistics — Unit 2, Lesson 2.9 — Student Packet Cover Sheet
\documentclass[10pt]{article}
\usepackage{apstats-article}
\usepackage{apstats-boxes}
\usepackage{ltablex}
\keepXColumns

\begin{document}

% Full-bleed navy banner
\begin{tikzpicture}[remember picture, overlay]
  \fill[navy] (current page.north west)
    rectangle ([yshift=-0.9in]current page.north east);
\end{tikzpicture}
\vspace{-0.8in}

\begin{center}
  {\color{white}\textbf{\LARGE AP Statistics}}\\[4pt]
  {\color{skymid}\large\textbf{Unit 2: Exploring Two-Variable Data}}\\[6pt]
  {\color{white}\textbf{\large Lesson 2.9 \quad Analyzing Departures from Linearity}}
\end{center}

\vspace{0.22in}
\namedateperiod
\vspace{0.18in}

% ── LEARNING TARGETS ─────────────────────────────────────────────────────────
\begin{learningtargetbox}
\small
\begin{enumerate}[leftmargin=1.5em, itemsep=5pt]
  \item \ldots{} identify when a linear model is \textbf{not appropriate} by examining a
        scatterplot and residual plot for curved patterns.
        \hfill\textcolor{slate}{\small(DAT-1.G)}
  \item \ldots{} apply a \textbf{logarithmic transformation} to $y$, to $x$, or to both,
        and explain which model (exponential or power) each choice produces.
        \hfill\textcolor{slate}{\small(DAT-1.H)}
  \item \ldots{} write the \textbf{regression equation in transformed form} and
        \textbf{back-transform} to make a prediction on the original scale.
        \hfill\textcolor{slate}{\small(DAT-1.H)}
  \item \ldots{} assess whether a transformation was successful by inspecting the
        \textbf{residual plot} of the transformed model.
        \hfill\textcolor{slate}{\small(DAT-1.G)}
\end{enumerate}
\end{learningtargetbox}

\vspace{0.14in}

% ── TABLE OF CONTENTS ────────────────────────────────────────────────────────
\begin{tocbox}
\small
\renewcommand{\arraystretch}{1.5}
\begin{tabularx}{\linewidth}{@{} c l X r @{}}
\rowcolor{sky}
\textbf{\#} & \textbf{Component} & \textbf{Description} & \textbf{Score} \\
\midrule
1 & Warm-Up       & Spiral review --- residual plots, regression, and logarithm properties  & \blank{1.2cm} \\
2 & Guided Notes  & Exponential and power transformations; back-transforming predictions     & \blank{1.2cm} \\
3 & Group Activity & Identifying, applying, and assessing transformations in context         & \blank{1.2cm} \\
4 & Exit Ticket   & Apply and back-transform a log model independently                      & \blank{1.2cm} \\
5 & Homework      & Multi-context transformation and prediction practice                     & \blank{1.2cm} \\
\midrule
  & \multicolumn{2}{r}{\textbf{Total}} & \blank{1.2cm} \\
\end{tabularx}
\end{tocbox}

\vspace{0.14in}

% ── KEEP IN MIND ─────────────────────────────────────────────────────────────
\begin{remindbox}
\small
When a scatterplot shows a curve, a linear model is not appropriate. Log transformations can
``straighten'' the relationship so regression applies. Here are the AP exam rules.

\vspace{6pt}
\begin{tabularx}{\linewidth}{@{} >{\centering\arraybackslash\columncolor{goldbg}}p{0.06\linewidth}
                                    X @{}}
\textbf{\large!} &
\textbf{Write the equation in transformed form, then back-transform for predictions.}\enspace
If you log-transformed $y$, write $\widehat{\log y} = a + bx$; then predict $y$ with
$\hat{y} = 10^{a+bx}$. Reporting the transformed value as the final answer loses credit. \\[6pt]
\textbf{\large!} &
\textbf{State which variable(s) were transformed.}\enspace
``I transformed $y$'' (exponential model) is different from ``I transformed both $x$ and $y$''
(power model). The AP exam requires you to specify. \\[6pt]
\textbf{\large!} &
\textbf{Verify with the residual plot.}\enspace
After transforming, check that the residual plot shows random scatter. A remaining pattern
means the transformation did not fully linearize the data. \\[6pt]
\textbf{\large!} &
\textbf{$r^2$ should increase after a good transformation.}\enspace
A higher $r^2$ for the transformed model (compared to the original) is evidence the
transformation improved the fit.
\end{tabularx}
\end{remindbox}

\end{document}
